\chapter{Keloid}

\section*{Definition}
A keloid is a benign, dermal fibroproliferative tumor extending beyond the margins of the original wound, distinguished from hypertrophic scars by growth beyond the injury boundary and lack of spontaneous regression. Incidence is highest in individuals of African, Asian, and Hispanic descent (15--20\%).

\section*{What Happens in Your Body}
Keloid formation results from aberrant wound healing with prolonged inflammatory phase, excessive fibroblast proliferation, and overproduction of collagen (primarily type III) and extracellular matrix components. Dysregulation of TGF-\ensuremath{\beta}1/Smad signaling, increased VEGF and PDGF, and decreased matrix metalloproteinases (MMPs) prevent normal collagen remodeling. Genetic predisposition with autosomal dominant inheritance and variable penetrance is recognized.

\section*{Causes}
\begin{itemize}
  \item \textbf{Cutaneous injury:} Surgical incisions, lacerations, burns, piercings, tattoos, vaccinations, acne, chickenpox scars
  \item \textbf{Sites of predilection:} Earlobes, sternum, shoulders, upper back, jawline, deltoid
  \item \textbf{Risk factors:} Darker skin types (Fitzpatrick IV--VI), family history, younger age (10--30 years), pregnancy, high-tension skin areas
  \item \textbf{Genetic syndromes:} Rubinstein--Taybi syndrome, Ehlers--Danlos syndrome, osteogenesis imperfecta
\end{itemize}

\section*{When to See a Doctor}
See your doctor if:
\begin{itemize}
  \item Scar continues to grow beyond wound margins for more than 3 months
  \item Pain, pruritus, or paresthesia at the scar site
  \item Functional impairment: restricted range of motion or cosmetic disfigurement affecting quality of life
  \item Rapid growth or new lesion without clear history of trauma (rule out dermatofibrosarcoma protuberans)
\end{itemize}

\section*{Related Symptoms}
\begin{itemize}
  \item Hypertrophic scar, dermatofibroma, dermatofibrosarcoma protuberans, scleroderma
\end{itemize}
\textbf{Important:} Keloids can recur after excision at rates of 50--100\% when surgery is used alone, making adjuvant therapy essential for optimal outcomes.

\section*{Self-Care / Home Management}
\begin{itemize}
  \item Apply silicone gel sheeting or silicone gel daily for at least 12 hours for 3--6 months on early scars
  \item Use compression earrings for earlobe keloids after piercing
  \item Sun protection on scars to prevent hyperpigmentation
\end{itemize}

\section*{How Your Doctor Will Evaluate You}
\begin{itemize}
  \item Clinical diagnosis based on growth beyond wound borders, absence of spontaneous regression, and typical location
  \item Skin biopsy if diagnosis uncertain or suspicion of malignancy (atypical features, rapid growth, ulceration)
  \item Evaluate for family history and prior keloid formation at multiple sites
\end{itemize}

\section*{Treatment Options}
\begin{itemize}
  \item First-line: Intralesional triamcinolone acetonide (10--40 mg/mL) every 4--6 weeks alone or with 5-FU
  \item Surgery: Intralesional excision with immediate adjuvant therapy (steroid injection, radiation, or pressure therapy)
  \item Radiation therapy: Superficial X-ray or brachytherapy within 24--48 hours post-excision
  \item Other: Cryotherapy, pulsed-dye laser, laser ablation, imiquimod cream, and verapamil injection
\end{itemize}

\section*{References}
\begin{thebibliography}{99}
\bibitem{keloid:ref1} Ogawa R. ``Keloid and hypertrophic scars are the result of chronic inflammation in the reticular dermis.'' \textit{International Journal of Molecular Sciences}, 2017;18(3):606.
\bibitem{keloid:ref2} Gauglitz GG, Korting HC, Pavicic T, et al. ``Hypertrophic scarring and keloids: pathomechanisms and current and emerging treatment strategies.'' \textit{Molecular Medicine}, 2011;17(1--2):113--125.
\bibitem{keloid:ref3} Mustoe TA, Cooter RD, Gold MH, et al. ``International clinical recommendations on scar management.'' \textit{Plastic and Reconstructive Surgery}, 2002;110(2):560--571.
\bibitem{keloid:ref4} Alster TS, Tanzi EL. ``Hypertrophic scars and keloids: etiology and management.'' \textit{American Journal of Clinical Dermatology}, 2003;4(4):235--243.
\bibitem{keloid:ref5} Monstrey S, Middelkoop E, Vranckx JJ, et al. ``Updated scar management practical guidelines.'' \textit{Burns}, 2014;40(6):1069--1082.
\bibitem{keloid:ref6} Tredget EE, Nedelec B, Scott PG, et al. ``Hypertrophic scars, keloids, and contractures.'' \textit{Surgical Clinics of North America}, 1997;77(3):701--730.

\end{thebibliography}
